\section{Excluding a tilted full-turn selector}\label{app:tilted-selector}

Throughout this appendix $U$ is a canonical full-turn maximum.
Reflect horizontally when necessary, and write
\[
 I=[-2C,2C],\quad J=[-C,C],\quad
 f_U(-C)=1-h,\quad f_U(C)=1,\quad h>0.
\]
The preliminary cuts and \cref{cor:full-three-tent-cut} give
$1/2<C<37/50$ and $0<h<17/50$.
Write $n=n_L(U)$, $n_\pm=n(\pm C)$ and $e_R=f_U(2C)$,
$e_L=f_U(-2C)$. The actual endpoint and wing identities are
\begin{equation}\label{eq:tilted-endpoints}
 \begin{gathered}
 e_R=\tfrac12+\tfrac h4+(3n_+-n_-)/4,\quad
 e_L=\tfrac12-\tfrac{3h}4+(3n_--n_+)/4,\\
 E=e_R+e_L=1-h/2+(n_++n_-)/2,\quad
 \max_Jn\le E/2,\quad 2P_L(U)=L_{\rm w}.
 \end{gathered}
\end{equation}
Let $F$ be the first support and $G$ the low-left-wing companion
support, with the uncharged middle facet removed, as in the tilted
support construction. Use $p,q,u,v$ for their velocities and regular
densities in \eqref{eq:variation-velocities}.

\subsection{A connected ordinary survivor}

\begin{lemma}\label{lem:tilted-box-confinement}
Every attached quadrant with positive apex height has its apex
abscissa in $(-C,C)$. The full niche is convex on each exterior
wing, and $n\le f_U$ on $I$. Consequently the survivor
$S=\{(x,y):x\in I,\ n(x)\le y\le f_U(x)\}$ is an ordinary sofa.
\end{lemma}
\begin{proof}
Lift the low wing by $h$ and make the middle roof one, obtaining
$\widehat U$. Write $s=\sin t$, $c=\cos t$ and let
$\widehat F,\widehat G$ be its two supports.
For a horizontal-middle cap the box and top endpoints imply
\[
 \widehat F\le2Cc+s,\quad\widehat G\ge Cs+c,
 \quad\widehat z_x\le C(2-3s^2)-c+s.
\]
Here $C(1-3s^2)\le c-s$ for $1/2\le C\le13/15$.
For $s^2\le1/3$ use $c\ge1-3s^2/5$; the margin is at least
$2s^2-s+2/15=2(s-1/4)^2+1/120$.
For $1/3\le s^2\le1/2$ the two sides have opposite signs.
For $s^2\ge1/2$ use
$2(s-c)=2(2s^2-1)/(s+c)\le4s^2-2\le3s^2-1$.
Reflection gives $-C\le\widehat z_x\le C$.

Put $\delta=\arctan(h/(2C))$. For $t\ge\delta$ the actual
second support is on the low wing, and
$F=\widehat F$, $G=\widehat G-hc$.
Thus $z_x=\widehat z_x+hsc>-C$ in the open angular interval.
To prove its upper bound it suffices to show
\begin{equation}\label{eq:tilted-corner-margin}
 c-s-C(1-3s^2)-hsc>0.
\end{equation}
For $s^2\le1/3$, its margin is at least
$[c(50-17s)-(37+50s-111s^2)]/50$.
Enlarge to $0\le s\le3/5$. The two compared terms are positive;
their squared difference is
\[
 Q(s)=1131-5400s+3503s^2+12800s^3-12610s^4.
\]
With $z=s-3/8\in[-3/8,9/40]$ this is
\[
 Q=49647/2048-(2091/64)z
       +z^2(116213/16-6115z-12610z^2).
\]
The bracket is concave and its endpoint values are $249061/32$
and $839849/160$, both above $5000$.
Hence $Q>24-33|z|+5000z^2\ge24-1089/20000>0$.
For $s^2\ge1/3$ put $C=1/2$, enlarge $h$ to $2/5$ and factor
its lower bound as
\[
 c\left[1-2s/5-\frac{c(3s+1)}{2(1+s)}\right].
\]
Squaring the nonnegative compared terms reduces positivity to
$R(s)=75-105s-139s^2+241s^3>0$.
For $z=2s-1\in[0,1]$,
$8R=143-253z+445z^2+241z^3\ge143-253^2/(4\cdot445)>0$.
This proves \eqref{eq:tilted-corner-margin}.
For $0<t<\delta$ the actual second support is $(C,1)$, and the
apex height is at most
$1-c-s+Csc=sc[C-2/(1+c+s)]<0$.
Such quadrants have no positive roof.

Every positive quadrant uses its descending first wall on the right
wing and its increasing second wall on the left wing.
The niche is therefore a maximum of affine functions and zero on
each wing, hence convex there; the box gives $n(\pm2C)=0$.
The endpoint box gives $n_\pm\le2/5$.
By \eqref{eq:tilted-endpoints},
$\max_J n\le7/10-h/4<1-h\le f_U$.
On a wing $f_U-n$ is concave and nonnegative at both endpoints,
so the same inequality holds throughout $I$.
Apply \cref{lem:ordinary-survivor}.
\end{proof}

In particular, with $G_0=22199/10000$, convexity of the wing niche
and the stationary chord bounds give
\begin{equation}\label{eq:tilted-ordinary-bound}
 P_L(U)+C(3-E-3h/2)\le |S|\le G_0,\qquad
 P_L(U)\ge\sqrt{C^2+(1-h/2-E/2)^2}.
\end{equation}
Indeed the middle cap area is $2C-Ch$, while the exterior niche
is at most $C(n_++n_-)/2=C(E-1+h/2)$.
The two wing chords have vectors $(C,1-e_R)$ and
$(C,1-h-e_L)$; their average length is at least the length of
their average vector.

\subsection{The forward excess ends before endpoint tangency}

\begin{lemma}\label{lem:tilted-late-unit}
If $2/3<C<37/50$, then $u(t)\le1$ for almost every $t\ge11/15$.
Moreover, if the top face is $[C,C+T]\times\{1\}$, then
\begin{equation}\label{eq:tilted-endpoint-leakage}
 \begin{gathered}
 n_+\le\frac{CT}{\sqrt{1-C^2}},\qquad
 J_F:=\int_0^L(u-1)_+\sin t\,\dd t<1/40,\\
 \int_C^{2C}n\le(T+J_F)n_+/2.
 \end{gathered}
\end{equation}
\end{lemma}
\begin{proof}
Write $D=3C$, $B=e_R-1+h$ and $\gamma_0=\arccos(1-h)$.
The endpoint box $H(C)=1-\sqrt{1-C^2}<1/3$ and the sharper
low-wing box $n_-\le(H(C)-h)_+$ imply
\[
 -1/2+5h/4-(1/3-h)_+/4\le B\le-1/4+5h/4,
 \qquad -7/12<B<7/40.
\]
Only the initial positive-$q$ component can have $u>1$.
As in the reflected calculation in \cref{app:partial-reflection},
later positive components have $q\le1/8$.
Before its first $p$-zero $\beta$ the same-sign law gives $u=0$.
If $q(\beta)\le1$ no excess occurs. Otherwise put
$\eta=q(\beta)-1>0$.
On its only possible excess episode, at elapsed time $r\le1$,
\[
 p\le-r,\quad q'\le-(1+r)/2,\quad
 (u(\beta+r)-1)_+\le(\eta-r/2-r^2/4)_+.
\]
The duration is at most $\ell(\eta)=\sqrt{1+4\eta}-1$.
All bounds below give $\ell<1$, as required.

If $\beta\le\gamma_0$, the initial-floor solution in
\eqref{eq:initial-floor-solution} gives
$(q(\beta)+1)^2=D^2+B^2-1>4$.
Thus $B<0$, since $(111/50)^2+(7/40)^2<5$.
Also $D\sin\beta=1+B\cos\beta\le1$, so
$\beta<\pi/6<8/15$.
Finally $(q(\beta)+1)^2<(111/50)^2+(7/12)^2-1<(207/100)^2$,
so $\eta<7/100$, $\ell<2/15$, and the episode ends before $2/3$.

If $\gamma_0<\beta$, then
$1+B(1-h)-D\sqrt{h(2-h)}>0$.
This forces $h<1/8$. Indeed its upper bound
\[
 V(h)=1+[-1/4+5h/4](1-h)-2\sqrt{h(2-h)}
\]
is convex on $[1/8,17/50]$ and negative at both endpoints:
$V(1/8)=235/256-\sqrt{15}/4<0$ and
$V(17/50)=1+231/2000-\sqrt{1411}/25<0$.
For $h<1/8$ the sharper interval is
$-7/12+3h/2\le B\le-1/4+5h/4<-3/32$.
The energy used in \cref{lem:tilted-initial-energy} gives
\[
 (q(\beta)+1)^2\le
 Q(D,B,h)=D^2+B^2+B(1-h)-D\sqrt{h(2-h)}.
\]
For $1/25\le h<1/8$, $Q$ increases with $D$, is convex in
$B$, and is convex in $h$ at either affine $B$ endpoint.
It suffices to evaluate $D=111/50$ and $h=1/25,1/8$.
At $h=1/25$ the larger endpoint is $B=-1/5$, giving
$Q=10387/2500<(51/25)^2$.
At $h=1/8$, $B^2+B(1-h)<0$ and
$\sqrt{h(2-h)}>15/32$ give a smaller value.
Hence $\eta<1/25$, $\ell<2/25$.
While $p,q>0$, $p(t)\le29/32-Dt+t^2/2$, negative at $t=3/5$;
therefore the episode ends before $17/25$.
For $0<h<1/25$, one has $-7/12<B<-1/5$, and
\[
 \beta\le D-\sqrt{D^2-8/5},\qquad
 \eta\le(\sqrt{D^2-4/25}-2)_+.
\]
The second inequality uses $B(1-h)-D\sqrt{h(2-h)}\le B$
and $B^2+B\le-4/25$.
On the range $D\ge\sqrt{104}/5$, the sum
\[
 W(D)=D-\sqrt{D^2-8/5}+\sqrt{4\sqrt{D^2-4/25}-7}-1
\]
increases: the first two terms have derivative above $-1/3$,
and the last term has derivative above $10/7$.
At $D=111/50$ the first part is below $2/5$ and the duration
below $1/3$, since $\sqrt{D^2-4/25}<219/100$.
Thus every episode ends before $11/15$.

The alternating Taylor bound gives
\[
 \cos(11/15)\ge6093080189/8201250000>37/50.
\]
At a positive endpoint-active angle at $x=C$, confinement gives
strict companion slack, so stationarity has $B_x(t)=C$.
The outer first-support abscissa is at most $2C$, hence
$\cos t\le C$ and $t\ge\arccos C>11/15$.
The backward support formula then gives
\[
 a_t(C)=T\cot t+\frac1{\sin t}
       \int_t^L\sin(r-t)(u(r)-1)\,\dd r
       \le T\cot t\le\frac{CT}{\sqrt{1-C^2}}.
\]
Also $(F(t)-1)/\cos t\le C+T+J_F$.
Thus the convex right-wing niche vanishes by $C+T+J_F$ and has
area at most $(T+J_F)n_+/2$.
In every excess case above $\eta<19/100$ and $\ell<1/3$.
The quadratic envelope has integral less than $11/324$,
and occurs below $11/15$. Consequently
$J_F<(11/15)(11/324)=121/4860<1/40$.
\end{proof}

\subsection{The low endpoint is dry and the top overhang is positive}

\begin{lemma}\label{lem:tilted-dry-endpoint}
For $2/3<C<37/50$, the top overhang satisfies $T>0$ and
\[
 n_-=0,\qquad T\le\sqrt{h(2-h)},\qquad
 T\le\sqrt{1-e_L^2}-C,\qquad T<1-C.
\]
The late-unit and leakage conclusions of
\cref{lem:tilted-late-unit} hold as well.
\end{lemma}
\begin{proof}
Let $x_1=1-2C$ and consider the second-wall envelope
$\Psi(x)=\sup_{0<t<L}(G(t)-1+x\sin t)/\cos t$ for $x<x_1$.
It is finite, convex and nondecreasing. Its endpoint-angle limits
are $-h$ and $-\infty$. It is negative for $x\le-2C$ and
approaches a positive value at $x_1$, since the outer endpoint
$(-2C,e_L)$ has $e_L>0$.
It has a unique zero $x_0<x_1$: a zero maximum is attained at
an interior angle with strictly positive affine slope, which gives
strictly negative values to its left and positive values to its right.
For $x\le x_1$, the pinned top point $(C,1)$ gives
\[
 F(t)-1-x\cos t\ge(3C-1)\cos t+\sin t-1
             \ge\cos t+\sin t-1>0.
\]
For $x_1<x<C$, angles near $L$ have first wall
$(C+T-x)(L-t)+o(L-t)>0$ and second wall tending to $+\infty$.
The extra high-wing companion walls are nonpositive on $J$.
Hence $\{x\in(-C,C):n(x)>0\}=(\max(-C,x_0),C)$.

In this geometry the finite source projection identity has an exact
limit. On a compact interval to the left of $x_0$, all surrogate
second walls have a uniform negative gap. The high-point second
walls also have a negative gap outside arbitrarily small endpoint
angle neighborhoods. The coarse source bounds make their remaining
mass and horizontal projection $O(\delta)+o_i(1)$ for an angle
cutoff $\delta$. On compact subsets of the positive interval,
uniform barrier convergence gives the opposite inclusion.
Let the approximation index tend to infinity and then $\delta\to0$.
The finite projection identity and source balance yield
\[
 2C-T=C-\max(-C,x_0),\qquad T=(C+x_0)_+.
\]
This argument controls shallow finite graph pieces explicitly.

Suppose $T=0$. Reflect the local support equations as
$P=-q(L-r)$, $Q=-p(L-r)$, $U=v(L-r)$, $V=u(L-r)$.
Before the removed middle normal they have initial data
$(P(0),Q(0))=(e_L,3C-1)$ and the same arm and shadow laws.
Beyond that normal $U=0$ directly.
The box and endpoint equations give $3/20<e_L<3/4$.
Let $D=3C$ and $e=e_L$.
At the first $P$-zero, the energy calculation without an impulse
gives $(Q+1)^2\le D^2-e(1-e)$ and a crossing time at most
$D-\sqrt{D^2-2e}$.
The same quadratic excess envelope then applies.
For $e\le2/3$, $e(1-e)\ge51/400$, so the excess height is
below $1/5$, its duration below $7/20$, and the crossing time
below $2-\sqrt{8/3}<3/8$.
For $2/3\le e<3/4$, $e(1-e)\ge3/16$, so these bounds are
$9/50$, $5/16$ and $2-\sqrt{5/2}<21/50$.
The sums are $29/40$ and $293/400$, both below $11/15$.
Later positive-$Q$ components have amplitude at most $1/8$.
Thus $v(t)\le1$ whenever $L-t\ge11/15$.

If $\Psi(-C)\ge0$, its maximum is attained at a proper angle
$t$ and has shifted companion abscissa $D_x(t)=-C$.
The outer supporting abscissa is $D_x-\sin t\ge-2C$,
so $\sin t\le C$ and $L-t\ge\arccos C>11/15$.
The entire prefix $[0,t]$ then has $v\le1$; its support formula
from $G(0)=1-h$, $G'(0)=C$ gives
\[
 \frac{G(t)-1-C\sin t}{\cos t}
 =1-h-\sec t+\frac1{\cos t}\int_0^t v(r)\sin(t-r)\,\dd r
 \le-h<0.
\]
This contradiction shows $\Psi(-C)<0$, hence $x_0>-C$,
contrary to $T=(C+x_0)_+=0$.
Therefore $T>0$, $x_0=-C+T$ and $n_-=0$.

The outer endpoint supplies a nonnegative second wall at
$x=-2C+\sqrt{1-e_L^2}$; the low middle point supplies one at
$x=-C+\sqrt{h(2-h)}$.
Comparing each with the unique zero $x_0$ gives the two square-root
bounds (if the test point is at least $x_1$, use $x_0<x_1$).
Finally $x_0<x_1$ implies $T<1-C$.
\end{proof}

\subsection{A sharper triangle payment for moderate tilt}

\begin{lemma}\label{lem:tilted-medium-width}
No canonical full-turn maximum has
$73/100\le C<37/50$ and $1/20\le h\le21/100$.
\end{lemma}
\begin{proof}
By \cref{lem:tilted-dry-endpoint}, put $z=n_+$ and $n_-=0$.
Then $e_R=1/2+h/4+3z/4$ and $e_L=1/2-3h/4-z/4$,
with $z<1/3$.
Use the lifted three-angle parameters in
\eqref{eq:three-tent-objective}. Put $r=\sqrt2$, $k=r-1$,
$c_8=\cos(\pi/8)$, $A=1-k$, $b=2-k$,
$d_0=3r-1$, $\gamma=1-r/2$, $t=2c_8-r$.
The outer deficits satisfy
$w_R\le1/2-h/4$ and $w_L<7/12+3h/4$.
The signed triangle and positive-floor calculation gives
\[
 \begin{gathered}
 P_L(U)\le G_h(a,u,v)+L_{\rm a}(h),\quad a=2C,
 \\
 L_{\rm a}=\gamma[(D-h/4)_+^2+(D+1/12+3h/4)_+^2],
 \quad D=1/2-t.
 \end{gathered}
\]
The geometric checks from \cref{lem:full-three-tent-score} remain
valid down to $C=73/100$: the diagonal apex is positive since
$C-h/2-k>0$; the endpoint sum gives
$m-h/2\ge a-1/2-h/4\ge363/400>179/200$; and the extra
gain triangles' outer zeros remain inside $J$ since $C<4/5$.
Thus the same clipped $G_h$ applies.

Set $a_*=73/50$, $B_0(a)=a-(a-k)^2/2$ and
$s_0(a)=d_0a/2-(4c_8-2)$.
The common support critical point at $a_*$ is
\[
 u_*=v_*=(a_*+k)/2+(b/9)[s_0(a_*)-Ah/4]+h/4.
\]
It is interior, unclipped and has positive $S$.
Indeed $140/99<r<99/70$, $923/1000<c_8<231/250$ imply
$2/3<s_0(a_*)<7/10$ and $s_0-Ah/4>0$.
Its right clipping quantity is
$-k/2+(b/9)s_0+(1/4-bA/36)h<0$ by the same bounds;
the left one is smaller by $h$.
The value and width derivative are
\[
 K(a_*,h)=B_0(a_*)-kh/2-h^2/8-\tfrac29[s_0(a_*)-Ah/4]^2,
\]
\[
 \partial_aG_h(a_*,u_*,v_*)
 =r-a_*-(2d_0/9)[s_0(a_*)-Ah/4]<0.
\]
Joint concavity and its supporting plane give $G_h(a,u,v)\le K$
for every $a\ge a_*$, including every clipping pattern.

Both positive parts in $L_{\rm a}$ are positive on this interval:
$1/2-2(231/250)+140/99<D<239/3500$, and the lower bound
exceeds $21/400$.
The linear coefficient of $K+L_{\rm a}$ is
$-k/2+As_0/9+\gamma(D+1/8)<-9/100$.
Its quadratic coefficient is $-1/8-A^2/72+5\gamma/8<1/16$.
Thus the derivative is below $-9/100+h/8<0$ and the largest
bound occurs at $h=1/20$.
There $s_0-A/80>43789/66000>53/80$ and
\[
 L_{\rm a}(1/20)<\tfrac3{10}[(7/125)^2+(19/100)^2]<3/250.
\]
As $B_0(a_*)=(123/50)k-529/5000$, we get
\[
 \begin{aligned}
 P_L(U)&<\frac{123}{50}\frac{29}{70}-\frac{529}{5000}
 -\frac{41}{3960}-\frac1{3200}\\
 &\hspace{5mm}-\frac29(53/80)^2+\frac3{250}
 =\frac{7550393}{9240000}<\frac{41}{50}<\frac\Mc2.
 \end{aligned}
\]
This contradicts the reference competitor.
\end{proof}

\subsection{The last small-height alternative}

\begin{lemma}\label{lem:tilted-small-height}
No canonical full-turn maximum has $2/3<C<37/50$ and
$0<h\le509/10000$.
\end{lemma}
\begin{proof}
Put $P=P_L(U)$ and $\mathcal S=S$, the ordinary survivor of
\cref{lem:tilted-box-confinement}. By \cref{thm:ordinary-sofa-bound},
\begin{equation}\label{eq:tilted-sh-2}
 |\mathcal S|\le G_0:=22199/10000.
\end{equation}
\paragraph{Exact endpoint and leakage information}

Write the top face as \([C,C+T]\times\{1\}\), and put

\[
\begin{gathered}
z=n(C),\qquad s=\sqrt{h(2-h)},\quad
K=\frac{C}{\sqrt{1-C^2}},\\
J_F=\int_0^{\pi/2}(u(t)-1)_+\sin t\,\dd t.
\end{gathered}
\]

\Cref{lem:tilted-dry-endpoint} and the actual pressure equations give

\[
T>0,\quad n(-C)=0,
\]

\begin{equation}\label{eq:tilted-sh-3}
e_R=\frac12+\frac h4+\frac{3z}{4},\qquad
e_L=\frac12-\frac{3h}{4}-\frac z4.
\end{equation}

The pinned low middle point and the endpoint estimate \eqref{eq:tilted-endpoint-leakage} imply

\begin{equation}\label{eq:tilted-sh-4}
T\le s,\qquad 0\le z\le KT,\qquad K<\frac{10}{9}.
\end{equation}

For completeness, the first inequality follows from the signed second
wall envelope at its zero \(x_{\rm zero}=-C+T\). The low middle point
\((-C,1-h)\) gives the lower wall
\(1-h-\sec t+T\tan t\). Its value at \(\sin t=T\) is
\(1-h-\sqrt{1-T^2}\), and cannot be positive at an envelope zero.
Thus \(T^2\le h(2-h)\). We have \(T\le C<1\), so this angle is
proper. The bound on K follows by squaring and using
\(81\cdot1369<100\cdot1131\).

The ordinary box estimate also gives

\begin{equation}\label{eq:tilted-sh-5}
z\le1-\sqrt{1-C^2}<\frac{33}{100}.
\end{equation}

The early-end proof bounds all possible first-source excess by one
initial envelope

\[
(u(\alpha+\tau)-1)_+
\le\left(\varepsilon-\frac\tau2-\frac{\tau^2}{4}\right)_+,
\qquad \varepsilon<\frac{19}{100}<\frac7{36}.
\]

Its duration is less than \(1/3\), and all excess occurs at angles
less than \(11/15\). Since the parabola with initial value \(7/36\)
vanishes at \(1/3\),

\[
\int(u-1)_+\,dt
<\int_0^{1/3}\left(\frac7{36}-\frac\tau2-\frac{\tau^2}{4}\right)d\tau
=\frac{11}{324}.
\]

Consequently

\begin{equation}\label{eq:tilted-sh-6}
J_F<\frac{11}{15}\frac{11}{324}
=\frac{121}{4860}<\frac1{40};
\qquad \frac1{40}-\frac{121}{4860}=\frac1{9720}.
\end{equation}

\Cref{lem:tilted-late-unit} and convexity of the actual niche on the right wing give

\begin{equation}\label{eq:tilted-sh-7}
N_{\rm out}\le\frac{(T+J_F)z}{2}.
\end{equation}

The left wing contributes no outside niche: its niche roof is convex
with zero values at both endpoints. The middle cap area is \(2C-Ch\).
Thus

\begin{equation}\label{eq:tilted-sh-8}
|\mathcal S|\ge P+C(2-h)-\frac{(T+J_F)z}{2}.
\end{equation}

\paragraph{A common tangent lower bound for the wing lengths}

Define the fixed constants

\[
C_0=\frac{\sqrt{17}}6,\quad
R_0=\frac{\sqrt{26}}6,\quad
S_0=2C_0+R_0,
\]

\begin{equation}\label{eq:tilted-sh-9}
\lambda=\sqrt{\frac{17}{26}},\qquad
m=\frac{3}{4\sqrt{26}},\qquad
a=\frac{1-\lambda}{2}.
\end{equation}

The two charged-wing vertical drops are

\[
r_R=1-e_R=\frac12-\frac h4-\frac{3z}{4},\qquad
r_L=1-h-e_L=\frac12-\frac h4+\frac z4.
\]

They are nonnegative by the actual cap geometry. The top overhang has
length T, while the nonhorizontal right wing has horizontal displacement
\(C-T\). The stationary length identity and the two chord bounds yield

\begin{equation}\label{eq:tilted-sh-10}
P\ge\frac12\left[T+
\sqrt{(C-T)^2+r_R^2}+\sqrt{C^2+r_L^2}\right].
\end{equation}

The tangent inequality for the Euclidean norm at \((C_0,1/2)\) is

\[
\sqrt{x^2+y^2}\ge\lambda x+\frac{3}{\sqrt{26}}y.
\]

Apply it to the two chords in \eqref{eq:tilted-sh-10}, and use
\(\lambda C_0+3/(2\sqrt{26})=R_0\). Combining with \eqref{eq:tilted-sh-8} gives

\begin{equation}\label{eq:tilted-sh-11}
|\mathcal S|\ge S_0+(2+\lambda)(C-C_0)
-Ch-m(h+z)+aT-\frac{(T+J_F)z}{2}.
\end{equation}

The following rational enclosures follow by squaring positive numbers:

\begin{equation}\label{eq:tilted-sh-12}
\frac{8086}{10000}<\lambda<\frac{8087}{10000},\qquad
m<\frac{1471}{10000},\qquad a>\frac{956}{10000}.
\end{equation}

Since \eqref{eq:tilted-sh-4} implies \(T\ge9z/10\), \eqref{eq:tilted-sh-6} and \eqref{eq:tilted-sh-12} imply

\[
-mz+aT-\frac{J_Fz}{2}\ge-\frac{1839}{25000}z.
\]

Here

\[
\frac{1471}{10000}+\frac1{80}
-\frac{956}{10000}\frac9{10}
=\frac{1839}{25000}.
\]

In particular the useful common lower bound is

\begin{equation}\label{eq:tilted-sh-13}
|\mathcal S|\ge S_0+(2+\lambda)(C-C_0)
-\frac{8871}{10000}h-\frac{1839}{25000}z-\frac{Tz}{2}.
\end{equation}

Finally

\begin{equation}\label{eq:tilted-sh-14}
S_0>\frac{2\cdot4123+5099}{6000}
=\frac{2669}{1200}>G_0,
\quad \frac{2669}{1200}-G_0=\frac{16}{3750}.
\end{equation}

The radical comparisons here are
\((4123/1000)^2<17\) and \((5099/1000)^2<26\).

\paragraph{Failure of the second initial-energy inequality is impossible}

Put

\[
d=3C,\qquad B=e_R-1+h=-\frac12+\frac{5h+3z}{4}.
\]

The second condition of \cref{lem:tilted-initial-energy} is

\[
Q:=d^2+B^2+B(1-h)-ds\le4.
\]

Suppose it fails. The exact square completion is

\[
Q=\left(3C-\frac s2\right)^2
+\frac9{16}(h+z)^2-\frac14,
\]

so

\begin{equation}\label{eq:tilted-sh-15}
\left(3C-\frac s2\right)^2
+\frac9{16}(h+z)^2>\frac{17}{4}.
\end{equation}

Write \(w=h+z\). the stated height range and \eqref{eq:tilted-sh-5} give \(w<2/3\), and
\(3C-s/2>0\). Rationalizing the difference of the two positive
square roots gives

\[
\frac{\sqrt{17}}2-
\sqrt{\frac{17}{4}-\frac9{16}w^2}
=\frac{9w^2/16}{\sqrt{17}/2+
\sqrt{17/4-9w^2/16}}
<\frac9{64}w^2.
\]

The denominator exceeds four because \(w<2/3\). Therefore \eqref{eq:tilted-sh-15}
implies

\begin{equation}\label{eq:tilted-sh-16}
C-C_0>\frac s6-\frac3{64}(h+z)^2.
\end{equation}

Insert \eqref{eq:tilted-sh-16} into \eqref{eq:tilted-sh-13}. Use the lower enclosure for \(\lambda\) on
the positive s term, its upper enclosure on the negative squared term,
and \eqref{eq:tilted-sh-4} on all occurrences of z and T. The result is

\begin{equation}\label{eq:tilted-sh-17}
|\mathcal S|-S_0\ge
\frac{57955}{150000}s
-\frac{8871}{10000}h-\frac59s^2
-\frac{84261}{640000}
\left(h+\frac{10}{9}s\right)^2.
\end{equation}

The linear coefficient is exactly

\[
\frac{2+8086/10000}{6}
-\left(\frac{1471}{10000}+\frac1{80}
-\frac{956}{10000}\frac9{10}\right)\frac{10}{9}
=\frac{11591}{30000}=\frac{57955}{150000}.
\]

Here is a complete rational positivity certificate for \eqref{eq:tilted-sh-17}. Set

\[
A=\frac{11591}{30000},\quad
D=\frac{8871}{10000},\quad
Q_0=\frac{84261}{640000},\quad
k=\frac{10000}{19491},\quad s_* =\frac{63}{200}.
\]

Since \(h\le509/10000\),

\[
h\le ks^2,\qquad 0<s<s_*.
\]

Indeed \(s^2=h(2-h)\ge(19491/10000)h\), and

\[
s_*^2-\frac{509}{10000}\frac{19491}{10000}
=\frac{1581}{100000000}>0.
\]

The right side of \eqref{eq:tilted-sh-17} is consequently at least

\[
s\left[A-\left(Dk+\frac59\right)s
-Q_0s\left(\frac{10}{9}+ks\right)^2\right].
\]

The bracket decreases for \(s\ge0\). Its value at \(s_*\) is
exactly

\begin{equation}\label{eq:tilted-sh-18}
A-\left(Dk+\frac59\right)s_*
-Q_0s_*\left(\frac{10}{9}+ks_*\right)^2
=\frac{11102231813}{13507522880000}>0.
\end{equation}

Thus \(|\mathcal S|>S_0>G_0\), contradicting \eqref{eq:tilted-sh-2}. The second energy condition cannot fail in the stated range.

\paragraph{Failure of the first initial-energy inequality is impossible}

Suppose instead that

\begin{equation}\label{eq:tilted-sh-19}
9C^2+B^2>5.
\end{equation}

Write \(\xi=(5h+3z)/4\), so \(B=-1/2+\xi\).
the stated height range and \eqref{eq:tilted-sh-5} imply \(0<\xi<1/2\), hence \(-1/2<B<0\).
It follows from \eqref{eq:tilted-sh-19} that

\begin{equation}\label{eq:tilted-sh-20}
C>\frac{\sqrt{19}}6>\frac{29}{40}.
\end{equation}

Also

\[
B^2>5-9\left(\frac{37}{50}\right)^2
=\frac{179}{2500}>\frac{16}{225}.
\]

The last gap is \(11/22500\). Since B is negative,

\begin{equation}\label{eq:tilted-sh-21}
B<-\frac4{15},\qquad \xi<\frac7{30}.
\end{equation}

These inequalities give a useful improved endpoint bound

\begin{equation}\label{eq:tilted-sh-22}
z<\frac{27}{100}.
\end{equation}

To verify it, if \(h\le1/40\), then
\(s<9/40\) and \eqref{eq:tilted-sh-4} gives \(z<1/4\).
If \(h\ge1/40\), then
\(\xi\ge1/32+3z/4\); \eqref{eq:tilted-sh-21} therefore gives
\(z<97/360<27/100\).

Use \eqref{eq:tilted-sh-20} together with \(C_0<11/16\),
\(2+\lambda>14/5\), and \eqref{eq:tilted-sh-14}. Since the width difference is positive,

\[
S_0+(2+\lambda)(C-C_0)
>\frac{2669}{1200}+\frac{14}{5}
\left(\frac{29}{40}-\frac{11}{16}\right)
=\frac{559}{240}.
\]

Finally \(T\le s<63/200\). Substituting this, the stated height range and \eqref{eq:tilted-sh-22} into
\eqref{eq:tilted-sh-13} gives

\begin{equation}\label{eq:tilted-sh-23}
\begin{aligned}
|\mathcal S|
&>\frac{559}{240}
-\frac{8871}{10000}\frac{509}{10000}
-\left(\frac{1839}{25000}+\frac{63}{400}\right)
\frac{27}{100}\\
&=\frac{666488123}{300000000}.
\end{aligned}
\end{equation}

Its exact margin over the ordinary bound is

\begin{equation}\label{eq:tilted-sh-24}
\frac{666488123}{300000000}-G_0
=\frac{518123}{300000000}>0.
\end{equation}

This contradicts \eqref{eq:tilted-sh-2} and excludes failure of the
first energy condition throughout the stated height range.


Both energy conditions in \cref{lem:tilted-initial-energy} therefore
hold. That lemma gives $u\le1$ on the whole first quarter, and
\cref{lem:tilted-first-unit-exclusion} excludes this tilted maximum.
\end{proof}
